\section{Objetivos del proyecto}
\label{cap:Proposal}

En esta sección se presentan los objetivos que guiarán el desarrollo de este trabajo. En correspondencia directa con la pregunta de investigación planteada en la Sección~\ref{sec:problem_statement}, se formula el objetivo general orientado a definir el alcance técnico y metodológico del proyecto. Posteriormente, se detallan los objetivos específicos que estructuran de forma secuencial y verificable las fases del desarrollo del método computacional, la preparación de los datos, el modelamiento en el espacio vectorial compartida, la validación híbrida y la transferencia de resultados.

\subsection{Objetivo general}

Desarrollar y evaluar un \textbf{método computacional} basado en procesamiento de lenguaje natural, representación semántica vectorial y algoritmos de agrupamiento no supervisado (\textit{clustering}) para analizar, agrupar y comparar cuantitativamente el grado de alineación curricular de las mallas curriculares y perfiles de egreso declarados de las carreras de Ingeniería Informática en Chile con respecto a los estándares de referencia internacional (ACM/IEEE CS2013 y CC2020) y a las competencias demandadas por el mercado laboral tecnológico nacional (estratificado entre el sector TI e industrias no tecnológicas), cuantificando las diferencias estructurales y brechas formativas mediante métricas de similitud semántica en un espacio latente común $\mathbb{R}^d$, y validando los agrupamientos obtenidos mediante un marco de evaluación híbrido (cuantitativo y cualitativo).

\subsection{Objetivos específicos}

\begin{enumerate}[label=\arabic*., nosep]
	\item Construir un corpus curricular y de mercado estructurado, estratificado y normalizado, con el propósito de consolidar una base de datos de texto plano depurada y comparable que sirva de insumo para el análisis semántico de los programas informáticos en Chile, mediante la recopilación curada de mallas curriculares (747 asignaturas) y perfiles de egreso declarados de 16 programas universitarios chilenos, la estructuración digital de las directrices internacionales ACM/IEEE (CS2013 y CC2020), la extracción y estratificación de 200 ofertas laborales de LinkedIn (divididas entre el sector tecnológico especializado y empresas no-TI), el curado manual de etiquetas de referencia (\textit{ground truth}) clasificando las 747 asignaturas en las 18 \textit{Knowledge Areas} de la ACM, y la aplicación de técnicas avanzadas de preprocesamiento NLP (remoción de \textit{stop words} académicas, normalización ortográfica y lematización). \textbf{Criterio de calidad}: Cobertura comprobada del 100\% de las mallas y asignaturas del corpus muestral (747 asignaturas en 16 programas) con una tasa de datos corruptos o nulos $< 1.0\%$ tras la etapa de normalización y filtrado de ruido.

	\item Generar representaciones vectoriales densas (\textit{embeddings}) de los objetos curriculares y laborales en un espacio latente común ($\mathbb{R}^d$), a fin de mapear las asignaturas de las mallas, los perfiles declarados, los estándares ACM y las vacantes de mercado en una misma geometría continua que represente su semántica contextual y habilite el cálculo matemático de similitudes, utilizando modelos avanzados basados en transformadores multilingües (\textit{Sentence-Transformers} SBERT/MPNet) en combinación con algoritmos de reducción de dimensionalidad (como UMAP o t-SNE) para su visualización espacial. \textbf{Criterio de calidad}: Conservación de la estructura topológica local del espacio vectorial verificada mediante una métrica de fidelidad de proyección (\textit{Trustworthiness} $\ge 0.85$).

	\item Identificar perfiles formativos específicos y medir la alineación semántica y las diferencias de mallas, con el fin de detectar patrones de especialización implícitos entre las distintas universidades y cuantificar con rigor matemático la correspondencia y las brechas curriculares frente a las directrices de referencia y las necesidades de la industria, mediante la aplicación coordinada de una estrategia tripartita de agrupamiento no supervisado (K-Means para particiones convexas con centroides explícitos, Agrupamiento Jerárquico para dendrogramas de parentesco, y DBSCAN para el aislamiento de asignaturas atípicas u \textit{outliers}), identificando patrones estructurales \textit{(e.g., concentración temprana de ciencias básicas vs. especialización precoz; profundidad de la línea de programación; presencia de ciberseguridad o inteligencia artificial en la malla)} y perfiles formativos \textit{(e.g., Malla orientada a Ingeniería de Software, Malla enfocada en Ciencia de Datos, Malla de Infraestructura y Redes, Malla de Computación Teórica)}, junto al cálculo de similitud del coseno ($\cos \theta$), tasa de cobertura ($\% \text{Cobertura}$) e índice de brecha semántica ($1 - \cos \theta$). \textbf{Criterio de calidad}: Obtención de una partición coherente con un Coeficiente de Silhouette $\ge 0.25$ y un Índice de Davies-Bouldin $< 1.50$.

	\item Validar la calidad, estabilidad y aplicabilidad práctica de los agrupamientos y brechas obtenidas, a objeto de asegurar la consistencia técnica del método computacional y la robustez de los hallazgos para su uso en la gestión académica, a través de un marco de evaluación híbrido que combine métricas cuantitativas intrínsecas, validación externa contrastada contra la taxonomía ACM utilizando el \textit{ground truth} de 747 asignaturas etiquetadas mediante el Índice Rand Ajustado ($\text{ARI} \ge 0.35$), Información Mutua Normalizada ($\text{NMI} \ge 0.40$) y Pureza ($\text{Purity} \ge 0.70$), evaluación de estabilidad mediante re-muestreo bootstrap ($\text{ARI}_{\text{promedio}} \ge 0.75$) y juicio cualitativo por un panel de 3 expertos independientes evaluado con un coeficiente de acuerdo inter-evaluador de Fleiss/Cohen $\kappa \ge 0.60$. \textbf{Criterio de calidad}: Cumplimiento simultáneo de los umbrales de validación intrínseca, externa, de estabilidad y de concordancia inter-experto estipulados.

	\item Operacionalizar el método computacional y sus resultados analíticos en una interfaz interactiva de apoyo a la toma de decisiones (\textit{Human-in-the-Loop}), para transferir el conocimiento generado a los gestores de las instituciones de educación superior y apoyar las deliberaciones informadas en los procesos de rediseño curricular continuo, mediante la implementación de un tablero de control (\textit{dashboard}) interactivo en React que integre mapas semánticos de asignaturas de las mallas, matrices de calor de competencias y gráficos dinámicos de brechas formativas. \textbf{Criterio de calidad}: Despliegue funcional de un tablero de control navegable y responsivo con una latencia media de respuesta $< 2.0$ segundos por consulta interactiva y usabilidad validada sin errores críticos de renderizado.
\end{enumerate}